\section{Sense}
\label{sec:axioms}

\subsection{Infinity and separation}

% Hopefully to include the leanest set of axioms here.
\begin{foundation}
    \begin{equation}
        \vdash \forall a, b: ((\forall c: (c \in a \lequiv c \in b)) \implies a = b)
        \tag{$\mathsf{EXT}$}
        \label{axiom:extensionality}
    \end{equation}
\end{foundation}

\begin{notation}
    \[a \notin b \coloneqq \neg(a \in b)\]
    \[\norneg{\norin{a}{b}}\]
\end{notation}

\begin{align*}
    \assumption[1]{(\forall c: c \notin a) \wedge (\forall c: c \notin b)}\\
    \byrule[1]{\forall c: c \notin a}{\ref{rule:conjunction-elimination}}\\
    \byrule[1]{d \notin a}{\ref{rule:universal-elimination}}\\
    \assumption[2]{d \in a}\\
    \byrule[2]{\bot}{\ref{rule:negation-elimination}}\\
    \byrule[2]{d \in b}{\ref{rule:deductive-explosion}}\\
    \byrule[1]{d \in a \implies d \in b}{\ref{rule:conditional-introduction}}\\
    \byrule[1]{\forall c: c \notin b}{\ref{rule:conjunction-elimination}}\\
    \byrule[1]{d \notin b}{\ref{rule:universal-elimination}}\\
    \assumption[2]{d \in b}\\
    \byrule[2]{\bot}{\ref{rule:negation-elimination}}\\
    \byrule[2]{d \in a}{\ref{rule:deductive-explosion}}\\
    \byrule[1]{d \in b \implies d \in a}{\ref{rule:conditional-introduction}}\\
    \byrule[1]{d \in a \lequiv d \in b}{\ref{rule:equivalence-introduction}}\\
    \byrule[1]{\forall d: (d \in a \lequiv d \in b)}{\ref{rule:universal-introduction}}\\
    \axiom[1]{\forall e, f: ((\forall d: (d \in e \lequiv d \in f)) \implies e = f)}{EXT}\\
    \byrule[1]{(\forall d: (d \in a \lequiv d \in b)) \implies a = b}{\ref{rule:universal-elimination}}\\
    \byrule[1]{a = b}{\ref{rule:conditional-elimination}}\\
    \byrule{((\forall c: c \notin a) \wedge (\forall c: c \notin b)) \implies a = b}{\ref{rule:conditional-introduction}}\\
    \byrule{\forall a, b: (((\forall c: c \notin a) \wedge (\forall c: c \notin b)) \implies a = b)}{\ref{rule:universal-introduction}}
\end{align*}
\begin{equation}
    \vdash \forall a, b: (((\forall c: c \notin a) \wedge (\forall c: c \notin b)) \implies a = b)
    \label{rule:empty-set-at-most-one}
\end{equation}
